% STATUS: reviewed
\chapter{Duomenų modelis}\label{ch:data}

Viskas, ką sistema žino, saugoma vienoje PostgreSQL duomenų bazėje su plėtiniu \code{pgvector}:
teisės aktų bazė ir jos vektoriai, pareigų registras, įmonės ir jų dokumentai, ataskaitos,
programos dokumentai, instrukcijos, užduočių eilė, kiekvieno modelio iškvietimo įrašas ir
veiksmų žurnalas. Prie jos jungiasi tik API (Django); asistento paslauga skaito ir rašo per API
(\cref{ch:arch}). Šiame skyriuje lentelės aprašomos taip, kaip jas apibrėžia modelių failai
(kiekvienos API programėlės \file{<app>/models.py}). Skyrius pradedamas taisyklėmis, kurių
laikosi kiekviena lentelė, ir dviem ryšių schemomis, toliau programėlė po programėlės
aptariamos lentelės, svarbiausi jų stulpeliai ir kiekvieno JSON stulpelio struktūra, o baigiama
tuo, ką ataskaita nusikopijuoja, kad vėliau skaitytųsi taip pat, ir tuo, kas --- modelis, kodas
ar žmogus --- rašo kuriuos stulpelius. Ką su lentelėmis daro apdorojimo grandinės, aprašyta
kiekvienai grandinei skirtame skyriuje.

\section{Taisyklės}\label{sec:data:conventions}

Kiekviena API programėlė savo lentelėms taiko tas pačias taisykles; \file{users/models.py} jas
išdėsto savo antraštiniame komentare, o kitos programėlės jų laikosi.

\paragraph{Raktai.} Pirminis raktas yra \code{TextField}, kuriame tekstu laikomas 4 versijos
UUID, sukuriamas Python kode įrašymo metu (\code{id=str(uuid.uuid4())}); duomenų bazė neturi
sekos ir nieko nepriskiria. Kadangi identifikatorius egzistuoja dar prieš įrašant eilutę,
etapas, kuriantis šimtus eilučių, jas įrašo vienu \code{bulk_create} ir vis tiek žino kiekvieną
identifikatorių --- taip savo išvadas įrašo atitikties spragų analizė, o veiksmų planas paskui
tuos identifikatorius išvardija \code{finding_ids}. Dvi lentelės vietoj to turi natūralųjį
raktą: \code{rate_limit_windows} --- ribotuvo raktą, \code{model_prices} --- modelio
identifikatorių. Viena lentelė identifikatorius gauna iš išorės: \code{chat_messages} eilutė
išlaiko identifikatorių, kurį žinutei suteikė pokalbio klientas; \file{chat/store.py} jį
priima tik tada, kai jis atitinka \code{MESSAGE_ID} (raidės, skaitmenys ir \code{_ . : -}, ne
daugiau kaip 128 simboliai) ir dar nepriklauso kitam pokalbiui.

\paragraph{Laikas.} Laiko momentas yra \code{DateTimeField} su laiko juostą žinančia UTC
reikšme (\file{settings.py}: \code{USE_TZ = True}, \code{TIME_ZONE = "UTC"}). Joks stulpelis
neužsipildo savaime: \code{auto_now} nenaudojamas, o kiekvienas įrašas perduoda
\code{utc_now()} reikšmę (\file{common/timestamps.py}), kurią šio failo antraštinis komentaras
vadina \enquote{vieninteliu laikrodžiu, kurį skaito kodas}. Perduodamas per API, momentas yra
ISO~8601 eilutė, besibaigianti \code{+00:00}. Diena be laiko --- data, kurią įsigalioja
redakcija, data, nuo kurios taikomas aktas, veiksmo terminas, leidimo galiojimas --- yra
\code{DateField}.

\paragraph{Išvardijimai.} Leistinos būsenos ar rūšies reikšmės yra konstantų rinkinys
(\en{tuple}) modelio faile (\code{REPORT_STATUSES}, \code{FINDING_STATUSES}, \dots), o
\code{CheckConstraint} modelio \code{Meta.constraints} priverčia duomenų bazę atmesti bet kokią
kitą reikšmę. Apribojimai pavadinti pagal lentelę ir tai, ką jie tikrina (\code{reports_status_check}). API rodiniai prieš įrašydami reikšmę
patikrina ją pagal tą patį rinkinį --- pavyzdžiui, įkeliant dokumentą pagal
\code{DOCUMENT_KINDS} ---, taigi apribojimas saugo nuo reikšmės, kurią suklysdamas pateiktų pats
kodas.


Keli tekstiniai stulpeliai apribojimo sąmoningai neturi, ir skaitytojas su jais susidurs
toliau: \code{work_products.kind} (\enquote{vėliau pridėtai rūšiai reikia modulio, ne
migracijos}), \code{report_acts.doc_type} (jame laikomi katalogo tipai ir dar
\code{internal} --- eilutė, prie kurios kabinami įmonės vidaus reikalavimai), išvadų ir
veiksmų \code{category} (ribojama tik lentelėje \code{obligations}),
\code{report_findings.source_kind} ir \code{report_findings.override_status} (rodinys jį
tikrina pagal \code{FINDING_STATUSES}), \code{prompt_versions.source},
\code{chat_messages.role} ir lentelės \code{llm_calls} \code{outcome} bei \code{purpose}. Du
\code{report_citations} stulpeliai yra arba tušti, arba viena iš sąrašo reikšmių: jų
apribojimai aiškiai leidžia \code{NULL}.

\paragraph{JSON stulpeliai.} Struktūrizuotos reikšmės, kurios skaitomos ir rašomos ištisai ---
akto sankcijų perskaitymas, išvados įrodymai, ataskaitos statistika ---, yra
\code{JSONField} stulpeliai, saugomi kaip \code{jsonb}. Jų struktūrą nustato juos rašantis
kodas, ne duomenų bazė, o raktai juose dažniausiai rašomi kupranugario stiliumi
(\code{maxFineEur}, \code{promptVersionId}), kaip juos pateikia API, --- išvardijimo reikšmė,
naudojama kaip raktas, išlaiko savo rašybą (\code{not_applicable} objekte \code{stats.findings})
---, o patys stulpeliai pavadinti mažosiomis raidėmis su pabraukimais. Juos skaitantis kodas sąmoningai
atlaidus: pavyzdžiui, \code{inventory_citations} faile \file{companies/inventory.py} trūkstamą
ar netaisyklingą sąrašą perskaito kaip nuorodų nebuvimą, niekada ne kaip klaidą.
\Cref{sec:data:corpus,sec:data:companies,sec:data:analysis,sec:data:programme,sec:data:ops}
pateikia kiekvieno iš jų struktūrą.

\paragraph{Vektoriai ir viso teksto paieška.} Vektorius yra \code{VectorField} su 1\,536
matmenimis (\code{EMBED_DIMENSIONS}, deklaruota \file{corpus/models.py} ir
\file{companies/models.py} kaip \enquote{vienas modelis}). Modelis nustatomas
\file{settings.py} (\code{AI_EMBED_MODEL}, numatytasis
\code{gateway:azure/text-embedding-3-large}), o asistento paslauga prašo jo būtent tiek
matmenų (\code{EMBED_DIMENSIONS} jos faile \file{routes/internal.js}); pakeitus modelį,
reikia vektorių perskaičiavimo užduoties, ne schemos pakeitimo (\cref{ch:models}). Vektorius
laiko keturi stulpeliai. Abu tie, kuriuose ieškoma, --- \code{legal_chunks.embedding} ir
\code{company_document_chunks.embedding} --- turi po HNSW indeksą kosinusiniam atstumui
(\code{m=16}, \code{ef_construction=64}, operatorių klasė \code{vector_cosine_ops}) ir šalia jo
GIN indeksą pagal \code{to_tsvector('simple', text)} mišriosios paieškos raktažodžių daliai.
Konfigūracija \code{simple} pasirinkta todėl, kad, \file{corpus/models.py} žodžiais,
\enquote{lietuvių kalbai Postgres neturi kamienavimo, linksniavimą perima vektorinė dalis}. Du
pareigos vektoriai (\code{embedding}, \code{embedding_en}) indekso neturi: jie yra užklausa,
lyginama su indeksuotomis įmonės dokumentų ištraukomis, o ne tai, kur ieškoma.

\paragraph{Kas vyksta trinant.} Kiekvienas išorinis raktas nurodo, kas nutinka eilutei, kai
ištrinama eilutė, į kurią jis rodo, ir šie pasirinkimai lemia, kas po trynimo išlieka
(\cref{tab:data:ondelete}). Dalis, kuri be visumos neturi prasmės, išnyksta kartu su ja
(\code{CASCADE}); naršymui laikoma nuoroda išvaloma, o eilutė išlieka (\code{SET_NULL});
savininko negalima ištrinti, kol jam kas nors priklauso (\code{PROTECT}); o įrašas, kuris turi
pergyventi tai, ką įvardija, identifikatorių laiko paprastu tekstu, visai be išorinio rakto.

\begin{xltabular}{\textwidth}{@{}L{22mm}YL{44mm}@{}}
\caption{Keturios trynimo taisyklės ir kur kiekviena taikoma}\label{tab:data:ondelete}\\
\toprule
\hd{Taisyklė} & \hd{Kur taikoma} & \hd{Pasekmė} \\
\midrule
\endfirsthead
\toprule
\hd{Taisyklė} & \hd{Kur taikoma} & \hd{Pasekmė} \\
\midrule
\endhead
\bottomrule
\endlastfoot
\code{CASCADE} & Akto redakcijos, struktūriniai vienetai ir fragmentai bei jo pareigos; redakcijos vienetai ir fragmentai; vieneto vaikiniai vienetai; įmonės leidimai, taikomi standartai, rolės, ataskaitos, veiksmų registras, programos dokumentai ir vidaus reikalavimai (\code{obligations.company}); dokumento fragmentai; ataskaitos aktai, išvados, veiksmai ir nuorodų eilutės; užduoties etapai; instrukcijos versijos; pokalbio žinutės; naudotojo sesijos ir pokalbiai; standarto taikymas, kai ištrinamas standarto aktas. & Ištrynus aktą iš teisės aktų bazės, pašalinamas visas jo pareigų registras; ištrynus įmonę, pašalinamos visos apie ją parengtos ataskaitos. \\
\code{SET_NULL} & Kiekvienas stulpelis, įvardijantis veikusį naudotoją (\code{created_by}, \code{reviewed_by}, \code{approved_by}, \dots, taip pat \code{jobs.created_by}, \code{llm_calls.user}, \code{audit_log.user}); kiekviena ataskaitos eilučių nuoroda į teisės aktų bazę ir į įmonės dokumentus; pareigos redakcija, nuostata ir šaltinio dokumentas; \code{legal_documents.current_version}, \code{prompts.active_version}; \code{company_documents.company}; \code{reports.job}, \code{work_products.job}; išvados akto eilutė (\code{report_findings.act}); pokalbio įmonė ir rolė; leidimo įrodymų dokumentas ir standarto taikymo sertifikatas; ataskaitos, kurias įvardija registro veiksmas. & Ataskaita išlieka ištrynus aktą, pareigą ar dokumentą, kurį ji skaitė: ji turi savo kopijas (\cref{sec:data:snapshots}) ir netenka tik nuorodos. \\
\code{PROTECT} & \code{companies.owner}, \code{company_documents.owner}. & Naudotojo, kuriam priklauso įmonė ar dokumentas, negalima ištrinti, kol jie jam priklauso. Paprastas paskyros trynimas ją išjungia (\code{users.active}), todėl veiksmų žurnalo eilutės ir nuosavybė lieka nepaliestos; tik administratoriaus galutinis ištrynimas (\route{DELETE}{/api/users/{id}?purge=true}, \code{purge_account} faile \file{users/purge.py}) paskyrą ištrina visam laikui: pirmiausia jos dokumentus, tada įmones, tada pačią paskyrą. \\
Be išorinio rakto & \code{job_steps.prompt_version_id}; \code{llm_calls.thread_id}, \code{.job_step_id}, \code{.prompt_version_id}; \code{obligations.prompt_version_id}; \code{chat_messages.parent_id}; \code{audit_log.target_id}; \code{company_actions.plan_action_id}; \code{report_actions.finding_ids} ir kiekvienas identifikatorius JSON stulpelyje. & Įrašas, kas kurį modelį iškvietė, už kokią kainą ir pagal kurią instrukcijos versiją, lieka išsamus ir tada, kai pokalbio, užduoties etapo ar versijos, kuriuos jis įvardija, nebėra. \\
\end{xltabular}

Du trynimai padaro daugiau nei vien taisyklės. Trinant įmonę, pirmiausia atsiejami jos
dokumentai (\code{company_documents.company} išvaloma, \file{companies/api/company_views.py}):
failai priklauso savo savininkui, ne įmonei. Ištrynus įmonės dokumentą, iš jo perskaityti
reikalavimai (šaltinio \code{internal} \code{obligations}) lieka su tuščiu
\code{company_document}, ir ataskaita jų nebeskaito, nes vidaus reikalavimų etapas ima tik tų
dokumentų, kuriuos jis skaito, reikalavimus (\file{analysis/internal.py}). Senas kelių lentelių
eilutes --- pasibaigusias sesijas, baigtas užduotis, sąnaudų įrašus --- išvalo kasdienė
priežiūros užduotis (\cref{ch:arch,ch:ops}).

\section{Ryšiai iš pirmo žvilgsnio}\label{sec:data:figures}

\Cref{fig:data:law} rodo teisės, įmonės ir ataskaitų lenteles ir tai, kaip jos susijusios;
\cref{fig:data:ops} --- lenteles, kuriomis sistema veikia. Ištisinė linija žymi dalį, kuri
ištrinama kartu su visuma, taškinė --- nuorodą, kuri išvaloma, kai išnyksta jos taikinys, o
brūkšninė --- identifikatorių, laikomą tekstu be išorinio rakto. Schemose stulpelių nėra; juos
pateikia tolesni poskyriai.

\begin{figure}[H]
\centering
\begin{tikzpicture}[x=1mm, y=1mm, font=\sffamily\scriptsize,
    tbl/.style={svc, minimum height=7mm, inner ysep=3pt, inner xsep=4pt},
    own/.style={draw=ink, line width=.7pt},
    ref/.style={draw=muted, line width=.6pt, densely dotted}]
  \begin{scope}[on background layer]
    \draw[lane] (0,8) rectangle (68,-54);
    \draw[lane] (76,8) rectangle (150,-48);
    \draw[lane] (0,-60) rectangle (150,-114);
  \end{scope}
  \node[lanelabel, anchor=south west] at (1,-53) {teisė};
  \node[lanelabel, anchor=south west] at (77,-47) {įmonė};
  \node[lanelabel, anchor=south west] at (1,-113) {ataskaitos ir tai, kas po jų};

  % the law
  \node[tbl] (ld) at (22,-4) {\code{legal_documents}};
  \node[tbl] (lv) at (22,-18) {\code{legal_document_versions}};
  \node[tbl] (lu) at (22,-32) {\code{legal_units}};
  \node[tbl] (lc) at (22,-46) {\code{legal_chunks}};
  \node[tbl] (ob) at (54,-30) {\code{obligations}};

  % the company
  \node[tbl] (co) at (134,-22) {\code{companies}};
  \node[tbl, align=center] (cx) at (132,-1) {\code{company_credentials}\\ \code{company_standards}\\ \code{company_roles}};
  \node[tbl] (cd) at (98,-22) {\code{company_documents}};
  \node[tbl] (cc) at (100,-38) {\code{company_document_chunks}};

  % reports
  \node[tbl] (re) at (80,-70) {\code{reports}};
  \node[tbl] (ra) at (22,-80) {\code{report_acts}};
  \node[tbl] (rf) at (52,-102) {\code{report_findings}};
  \node[tbl] (rc) at (106,-86) {\code{report_citations}};
  \node[tbl] (rx) at (100,-106) {\code{report_actions}};
  \node[tbl] (ca) at (128,-70) {\code{company_actions}};
  \node[tbl] (wp) at (128,-104) {\code{work_products}};

  % the law
  \draw[own] (ld) -- node[lbl, right] {redakcijos} (lv);
  \draw[ref] ([xshift=-6mm]ld.south) -- node[lbl, left] {dabartinė} ([xshift=-6mm]ld.south |- lv.north);
  \draw[own] (lv) -- node[lbl] {vienetai} (lu);
  \draw[own] (lu) -- node[lbl] {fragmentai} (lc);
  \draw[own] ([yshift=1.2mm]lu.west) -- ++(-3,0) |- node[lbl, pos=.25, left] {tėvinis} ([yshift=-1.2mm]lu.west);
  \draw[own] (ld.east) -| node[lbl, pos=.3] {aktas} (ob.north);
  \draw[ref] (lu.east) -- node[lbl, pos=.5, below] {nuostata} (ob.west);

  % the company
  \draw[own] (co) -- node[lbl] {dokumentai} (cd);
  \draw[own] (cd) -- node[lbl] {fragmentai} (cc);
  \draw[own] (co) -- (cx);
  \draw[ref] (ob.east) -- node[lbl, pos=.45] {išskirta iš} (cd.west);
  \draw[own] (cx.north) -- ++(0,4) -| node[lbl, pos=.25] {įmonės taikomas standartas} (ld.north);

  % reports
  \draw[own] (co.south west) -- ++(0,-28) -| node[lbl, pos=.2] {ataskaitos} (re.north);
  \draw[own] (co.south) -- node[lbl, pos=.55] {registras} (co.south |- ca.north);
  \draw[own] (co.east) -- ++(4,0) |- node[lbl, pos=.3] {programa} (wp.east);
  \draw[own] (re) -- node[lbl, pos=.4] {aktai} (ra);
  \draw[own] (re) -- node[lbl, pos=.45] {išvados} (rf);
  \draw[own] (re.south east) -- node[lbl, pos=.55, right] {nuorodų patikra} (rc.north west);
  \draw[own] (re.south) -- node[lbl, pos=.6] {planas} (rx.north west);
  \draw[ref] (ra) -- node[lbl] {aktas} (rf);
  \draw[ref] (rf.north) -- node[lbl, pos=.62] {pareiga} (ob.south);
  \draw[ref] (ca.west) -- node[lbl, pos=.4] {ataskaita} (ca.west -| re.east);
\end{tikzpicture}
\caption{Teisės ir įmonės pusės. Ištisinė linija: dalis, ištrinama kartu su visuma; taškinė:
  nuoroda, išvaloma ištrynus jos taikinį. Savo išvadas \code{report_actions} įvardija
  ne išoriniu raktu, o identifikatoriais JSON sąraše (\code{finding_ids}).}\label{fig:data:law}
\end{figure}

\begin{figure}[tbp]
\centering
\begin{tikzpicture}[x=1mm, y=1mm, font=\sffamily\scriptsize,
    tbl/.style={svc, minimum height=7mm, inner ysep=3pt, inner xsep=4pt},
    own/.style={draw=ink, line width=.7pt},
    ref/.style={draw=muted, line width=.6pt, densely dotted},
    txt/.style={draw=mid, line width=.6pt, dashed}]
  \node[tbl] (us) at (75,0) {\code{users}};
  \node[tbl] (se) at (20,0) {\code{sessions}};
  \node[tbl] (au) at (130,0) {\code{audit_log}};
  \node[tbl] (ct) at (26,-20) {\code{chat_threads}};
  \node[tbl] (cm) at (40,-40) {\code{chat_messages}};
  \node[tbl] (jo) at (75,-20) {\code{jobs}};
  \node[tbl] (js) at (75,-40) {\code{job_steps}};
  \node[tbl] (pr) at (130,-20) {\code{prompts}};
  \node[tbl] (pv) at (130,-40) {\code{prompt_versions}};
  \node[tbl] (lc) at (75,-60) {\code{llm_calls}};
  \node[tbl] (mp) at (130,-60) {\code{model_prices}};
  \node[ext] (cx) at (20,-60) {įmonės, rolės};
  \node[ext, align=center] (rw) at (51,-20) {ataskaitos,\\ programos dok.};
  \node[tbl, align=center] (st) at (124,-76) {\code{mcp_servers}\quad\code{rate_limit_windows}};

  \draw[own] (us) -- node[lbl] {sesijos} (se);
  \draw[ref] (us) -- node[lbl] {naudotojas} (au);
  \draw[own] (us.south west) -- node[lbl, pos=.45] {pokalbiai} (ct.north east);
  \draw[ref] (us) -- node[lbl] {created\_by} (jo);
  \draw[ref] (us.south east) -- node[lbl, pos=.55] {created\_by} (pv.north west);
  \draw[own] (ct) -- node[lbl] {žinutės} (cm);
  \draw[ref] (ct.south) -- node[lbl, pos=.5, left] {įmonė, rolė} (cx.north);
  \draw[own] (jo) -- node[lbl] {etapai} (js);
  \draw[own] (pr) -- node[lbl, left] {versijos} (pv);
  \draw[ref] ([xshift=6mm]pr.south) -- node[lbl, right] {aktyvi} ([xshift=6mm]pr.south |- pv.north);
  \draw[txt] (js) -- node[lbl, pos=.62] {prompt\_version\_id} (pv);
  \draw[txt] (lc) -- node[lbl] {job\_step\_id} (js);
  \draw[txt] (lc.north east) -- node[lbl, pos=.45] {prompt\_version\_id} (pv.south west);
  \draw[txt] (lc.west) -- node[lbl, pos=.4] {thread\_id} (ct.south east);
  \draw[txt] (lc) -- node[lbl] {kaina pagal modelio pavadinimą} (mp);
  \draw[ref] (rw.east) -- node[lbl, above] {užduotis} (jo.west);
  \draw[ref] ([xshift=3mm]us.south) .. controls (92,-12) and (92,-48) .. node[lbl, pos=.3] {naudotojas} (lc.north east);
  \node[note, anchor=north, align=center] at (st.south) {ryšių neturi, bet \code{mcp_servers.created_by}\\ rodo į \code{users} (išvaloma, kai naudotojo nebėra)};
\end{tikzpicture}
\caption{Veikimo lentelės. Ištisinė linija: dalis, ištrinama kartu su visuma; taškinė: nuoroda,
  išvaloma ištrynus jos taikinį; brūkšninė: identifikatorius, laikomas tekstu be išorinio rakto.
  Du brūkšniniai stačiakampiai žymi \cref{fig:data:law} lenteles.}\label{fig:data:ops}
\end{figure}

\section{Teisės aktų bazė ir pareigų registras}\label{sec:data:corpus}

Programėlė \code{corpus} laiko teisės aktus, jų redakcijas, jų tekstą, suskaidytą į
struktūrinius vienetus ir fragmentus, ir iš jų išskirtas pareigas. Aktas yra
\code{legal_documents} eilutė; kiekviena atsisiųsta ar įkelta redakcija yra
\code{legal_document_versions} eilutė; redakcijos struktūra yra \code{legal_units} medis; o
ištraukos, kuriose ieškoma, yra \code{legal_chunks}. Kaip tekstas ten patenka --- šaltiniai,
struktūros atpažinimas, skaidymas į fragmentus ---, aprašyta \cref{ch:corpus}; kaip pareigos
išskiriamos ir vertinamos --- \cref{ch:obligations}. 2026 m. lapkričio 13 d. lentelėse buvo 78
aktai, 71 redakcija, 28\,886 vienetai, 4\,059 fragmentai ir 5\,960 pareigų.

\begin{xltabular}{\textwidth}{@{}L{38mm}Y@{}}
\caption{\code{legal_documents}: vienas katalogo aktas ar standartas}\label{tab:data:ld}\\
\toprule
\hd{Stulpelis} & \hd{Reikšmė} \\
\midrule
\endfirsthead
\toprule
\hd{Stulpelis} & \hd{Reikšmė} \\
\midrule
\endhead
\bottomrule
\endlastfoot
\code{slug} & Unikalus; raktas, kuriuo aktą įvardija katalogas (\file{corpus/catalog.yaml}). \\
\code{title}, \code{title_en}, \code{short_title} & Oficialus pavadinimas, jo vertimas į anglų kalbą ir trumpasis pavadinimas, kurį spausdina nuorodos (\enquote{MiCA}; standarto atveju --- jo žymuo, \enquote{ISO/IEC 27001:2022}). \\
\code{jurisdiction} & \code{EU}, \code{LT} arba \code{INTL} (tarptautinis standartas). \\
\code{doc_type} & \code{regulation}, \code{directive}, \code{law}, \code{resolution}, \code{order}, \code{guideline}, \code{licensing}, \code{standard}, \code{other}. \\
\code{source_kind}, \code{external_id}, \code{source_url} & Iš kur gaunamas tekstas: \code{cellar}, \code{etar}, \code{eseimas} (senasis \code{etar} pavadinimas), \code{lb}, \code{url}, \code{upload}; identifikatorius šaltinyje (CELEX numeris, e-TAR identifikatorius) ir oficiali nuoroda. \\
\code{language} & Teksto kalba, numatytoji \code{lt}. \\
\code{domains}, \code{applies_to} & JSON sąrašai: akto temų žymos ir įstaigų tipai, kuriuos jis paprastai reglamentuoja (\code{casp}, \code{emi}, \dots{} arba \code{all}), --- pagal juos taikytinumo atranka filtruoja prieš modeliui perskaitant aktą. \\
\code{applicability_notes} & Kataloge parašyta pastaba tai atrankai. \\
\code{status} & \code{in_force}, \code{not_yet_applicable}, \code{repealed}. \\
\code{enabled} & Neigiama reikšmė (\en{False}) paslepia aktą nuo paieškos ir analizės, neištrindama jo redakcijų. \\
\code{current_version} & Naudojama redakcija; paieška visada skaito tik dabartinės redakcijos fragmentus. \\
\code{last_checked_at}, \code{last_error} & Paskutinis atnaujinimas ir jo klaida. \\
\end{xltabular}

\begin{xltabular}{\textwidth}{@{}L{38mm}Y@{}}
\caption{\code{legal_document_versions}: viena nekeičiama akto redakcija}\label{tab:data:lv}\\
\toprule
\hd{Stulpelis} & \hd{Reikšmė} \\
\midrule
\endfirsthead
\toprule
\hd{Stulpelis} & \hd{Reikšmė} \\
\midrule
\endhead
\bottomrule
\endlastfoot
\code{document}, \code{version_no} & Aktas ir redakcijos numeris; kartu unikalūs. \\
\code{content_hash}, \code{label} & Neapdorotų baitų maiša ir tai, kaip redakciją vadina šaltinis (konsoliduotas CELEX numeris, pvz., \code{02023R1114-20240109}, e-TAR redakcijos identifikatorius). Atnaujinimas, radęs tą pačią maišą ar žymą, nieko nesukuria, nebent jis priverstinis (užduoties \code{force}, \file{corpus/ingest.py}). \\
\code{file_path}, \code{content_type}, \code{source_url} & Neapdoroti baitai diske, įkėlimų aplanke, kad pakartotiniam struktūros atpažinimui nereikėtų iš naujo atsisiųsti. \\
\code{effective_date} & Data, kurią redakcija įsigalioja, jei žinoma. \\
\code{parser}, \code{text} & Struktūros atpažinimo modulis, kuris ją perskaitė (tuščias, kai tekstas tiesiog ištrauktas), ir ištrauktas grynas tekstas. \\
\code{char_count}, \code{unit_count}, \code{article_count}, \code{chunk_count}, \code{embed_model} & Kiekiai ir jos fragmentų vektorizavimo modelis. \\
\code{status}, \code{error} & \code{pending}, \code{parsing}, \code{embedding}, \code{ready}, \code{failed}. \\
\code{obligations_extracted_at} & Kada paskutinį kartą baigtas pareigų išskyrimas; tuščia reiškia, kad niekada, ir tada ataskaita pirmiausia išskiria pareigas. \\
\code{scope}, \code{sanctions}, \code{changes} & JSON: modelio perskaitytos akto taikymo srities ir sankcijų nuostatos bei tai, ką redakcija pakeitė, palyginti su ankstesne (žr. toliau). Tuščia reiškia, kad neperskaityta arba neužfiksuota. \\
\end{xltabular}

Trys redakcijos JSON stulpeliai yra tokios struktūros:

\begin{codeblock}
scope     = {types: [institution types, or "all"], reasoning,
             promptVersionId, model, readAt}
sanctions = {tiers: [{path, where, condition, documentId, infringedAct,
                      articles: [[from, to]], all, addressee, persons,
                      maxFineEur, maxFinePercent, benefitMultiple, measures}],
             note, provisions: [paths read], promptVersionId, model, readAt}
changes   = {from: previous version id, changed, added, removed: [paths]}
\end{codeblock}

\code{scope} rašo \file{corpus/scope.py}; jis leidžia taikytinumo atrankai palikti aktą,
kurį katalogo \code{applies_to} atmestų. Stulpelyje \code{sanctions} (\file{corpus/sanctions.py})
kiekviena pakopa yra maksimali bauda ar priemonė, kurią nustato viena nuostata: jos
\code{documentId} yra aktas, kurį ji liečia (šis, kitas teisės aktų bazės aktas arba tuščia,
o \code{infringedAct} tada įvardija aktą, kurio bazėje nėra), \code{articles} --- straipsnių
intervalai, kuriuos ji apima, kiekvienas nuo pirmojo iki paskutiniojo straipsnio pagal akto
jiems duotus numerius --- iki 2026~m. lapkričio 19~d. atliktame skaityme vien skaičiai ---
(\code{all}, kai apima visus straipsnius), \code{persons} --- asmenų rūšys, kurias pakopa
saisto, kai jos nuostata saisto tik dalį tų, kuriuos reguliuoja aktas, --- išlaikomos
straipsnių nenurodančioje pakopoje ir straipsnius nurodančioje pakopoje, šalia kurios kita
tų pačių straipsnių pakopa saisto kitus asmenis (kitu atveju tuščia, o iki tos dienos
atliktame skaityme šio lauko nėra), \code{addressee} ---
viena iš \code{legal}, \code{natural}, \code{any}, o \code{measures} --- kodai iš
\code{withdrawal}, \code{restriction}, \code{cease}, \code{public_statement},
\code{management_ban}, \code{disgorgement}, \code{periodic_penalty}, \code{criminal},
\code{other}; išlaikoma ne daugiau kaip \code{MAX_TIERS = 24}. \code{changes} apskaičiuojamas be
modelio, kai aktyvinama nauja redakcija (\code{version_changes} faile \file{corpus/ingest.py}),
straipsnis po straipsnio --- arba, jei akte straipsnių nėra, pagal aukščiausio lygio punktus
---, lyginant tekstus ir nepaisant tarpų.

\paragraph{Vienetai ir fragmentai.} Redakcijos tekstas saugomas dar dukart: kaip struktūrinių
vienetų medis, kurį įvardija nuoroda, ir kaip fragmentai, kuriuose ieškoma
(\cref{tab:data:lu}).

\begin{xltabular}{\textwidth}{@{}L{38mm}Y@{}}
\caption{\code{legal_units} ir \code{legal_chunks}}\label{tab:data:lu}\\
\toprule
\hd{Stulpelis} & \hd{Reikšmė} \\
\midrule
\endfirsthead
\toprule
\hd{Stulpelis} & \hd{Reikšmė} \\
\midrule
\endhead
\bottomrule
\endlastfoot
\multicolumn{2}{@{}l}{\textit{legal\_units}} \\
\code{version}, \code{parent}, \code{ordinal}, \code{depth} & Mazgo vieta redakcijos medyje, dokumento tvarka (indeksuota pagal redakciją ir eilės numerį). \\
\code{unit_type} & \code{part}, \code{title}, \code{chapter}, \code{section}, \code{subsection}, \code{article}, \code{paragraph}, \code{point}, \code{subpoint}, \code{annex}, \code{recital}, \code{preamble}, \code{text}, \code{clause} (numeruotas standarto skirsnis, \enquote{6.1.2}). \\
\code{number}, \code{heading} & Jo paties numeris ir antraštė; redakcija, tipas ir numeris kartu indeksuoti, kad būtų galima rasti cituojamą straipsnį. \\
\code{path} & Nuoroda, kurią jis spausdina: \enquote{59 str. 1 d.}, \enquote{Art. 59(1)}. \\
\code{text} & Jo \emph{visas} tekstas su vaikiniais vienetais, kad vien straipsnio eilutė atsakytų į \enquote{pacituok 59 straipsnį}. \\
\midrule
\multicolumn{2}{@{}l}{\textit{legal\_chunks}} \\
\code{version}, \code{unit}, \code{ordinal} & Paieškos vienetas --- straipsnis arba ilgo straipsnio dalių seka --- ir jo vieta (indeksuota pagal redakciją ir eilės numerį). \\
\code{document} & Akto kopija, kad paieška galėtų filtruoti pagal aktą be jungimo (indeksuota). \\
\code{path}, \code{heading}, \code{text}, \code{char_count} & Tai, ką rodo paieškos rezultatas; GIN indeksas apima \code{text}. \\
\code{embedding} & Po HNSW indeksu; tuščias, kol neatliktas vektorizavimo etapas, ir paieška tokias eilutes praleidžia. \\
\end{xltabular}

\begin{xltabular}{\textwidth}{@{}L{38mm}Y@{}}
\caption{\code{obligations}: viena pareiga arba vienas reikalavimas}\label{tab:data:ob}\\
\toprule
\hd{Stulpelis} & \hd{Reikšmė} \\
\midrule
\endfirsthead
\toprule
\hd{Stulpelis} & \hd{Reikšmė} \\
\midrule
\endhead
\bottomrule
\endlastfoot
\code{source_kind} & \code{legal} (iš teisės akto), \code{internal} (įmonės vidaus reikalavimas, perskaitytas iš jos politikos), \code{standard} (standarto reikalavimas, išskirtas arba įvestas administratoriaus). \\
\code{document}, \code{version}, \code{unit} & Aktas, redakcija, iš kurios perskaityta, ir pirmoji cituojama nuostata (\code{legal} ir \code{standard} atveju). \\
\code{company}, \code{company_document} & \code{internal} atveju: įmonė ir politika, iš kurios perskaityta. \\
\code{title}, \code{title_en}, \code{text}, \code{text_en} & Pareiga lietuvių ir anglų kalbomis. \\
\code{addressee}, \code{conditions}, \code{deadline}, \code{evidence} & Kam ji privaloma, kokiomis sąlygomis, iki kada ir kas įrodytų atitiktį. \\
\code{applies_to} & JSON sąrašas įstaigų tipų (arba \code{all}), kuriems pareiga privaloma; ataskaita palieka pareigas, tinkančias įmonei. \\
\code{category} & Viena iš dvidešimties: \code{licensing}, \code{registration}, \code{governance}, \code{policy}, \code{procedure}, \code{risk_management}, \code{ict_security}, \code{outsourcing}, \code{incident_reporting}, \code{reporting}, \code{disclosure}, \code{record_keeping}, \code{customer_protection}, \code{aml_kyc}, \code{capital}, \code{safeguarding}, \code{training}, \code{data_protection}, \code{ai_governance}, \code{other}. \\
\code{severity} & \code{critical}, \code{high}, \code{medium}, \code{low}. \\
\code{citations} & JSON: nuostatos, kuriomis ji grindžiama (žr. toliau). \\
\code{status} & \code{proposed} (iš modelio), \code{approved}, \code{rejected}. Atmesta eilutė niekada nepatenka į ataskaitą; patvirtinta eilutė išlieka pakartotinio išskyrimo metu. \\
\code{prompt_version_id}, \code{model} & Ją išskyrusi instrukcijos versija ir modelis, tekstu. \\
\code{verification} & JSON: vertintojo verdiktas (žr. toliau); tuščias, jei niekada nevertinta arba po to redaguota. Kai nė vienos cituojamos nuostatos dabartinėje akto redakcijoje nėra, kodas be modelio iškvietimo įrašo \code{unsupported} (\file{corpus/verify.py}). \\
\code{embedding}, \code{embedding_en} & Pavadinimo ir teksto vektoriai kiekviena kalba. \\
\code{notes}, \code{reviewed_by}, \code{reviewed_at} & Peržiūrėtojo pastaba ir kas pareigą patvirtino ar atmetė. \\
\end{xltabular}

\begin{codeblock}
citations (legal, standard) = [{unitId, path}]
citations (internal)        = [{path: the section, documentId, documentTitle,
                                quote, start}]
verification = {verdict: supported | partly | unsupported,
                issue: none | duty | addressee | condition | deadline | merged,
                note, promptVersionId, model, checkedAt}
\end{codeblock}

Vidaus reikalavimo \code{start} yra simbolis, nuo kurio jo citata prasideda politikos tekste,
arba tuščia, kai tekste ji pažodžiui neužrašyta. Patikra \enquote{nustato žmogaus peržiūros
eiliškumą ir nieko nesprendžia} (\file{corpus/models.py}). Registrą aptarnauja du indeksai:
pagal aktą ir būseną bei pagal įmonę ir būseną.

\section{Įmonės ir jų dokumentai}\label{sec:data:companies}

Programėlė \code{companies} laiko profilį, kuriuo remiasi ataskaita, tai, ką įmonė turi ir
taiko, jos roles ir jos dokumentus su jų tekstu, inventoriumi ir ištraukomis
(\cref{ch:company}). Įmonė priklauso savo savininkui; dokumentas turi savo savininką, nes failą
galima įmesti į pokalbį dar prieš atsirandant bet kokiai įmonei ir priskirti vėliau.

\begin{xltabular}{\textwidth}{@{}L{38mm}Y@{}}
\caption{\code{companies}: profilis}\label{tab:data:co}\\
\toprule
\hd{Stulpelis} & \hd{Reikšmė} \\
\midrule
\endfirsthead
\toprule
\hd{Stulpelis} & \hd{Reikšmė} \\
\midrule
\endhead
\bottomrule
\endlastfoot
\code{owner} & Naudotojas, kuriam ji priklauso (\code{PROTECT}). \\
\code{name}, \code{registration_code}, \code{country}, \code{legal_form} & Tapatybė; numatytoji šalis --- \code{LT}. \\
\code{nace_codes}, \code{sector}, \code{activities} & EVRK (NACE) kodų JSON sąrašas; laisvas tekstas. \\
\code{size_category}, \code{employees}, \code{annual_turnover_eur} & Dydis: \code{micro}, \code{small}, \code{medium}, \code{large}, \code{unknown}. \\
\code{institution_type} & Pagal ką pirmiausia filtruoja taikytinumo atranka: \code{bank}, \code{credit_union}, \code{payment_institution}, \code{emi}, \code{casp}, \code{crowdfunding}, \code{consumer_credit}, \code{real_estate_credit}, \code{currency_exchange}, \code{investment_firm}, \code{insurance}, \code{other_financial}, \code{none}, \code{unknown}. \\
\code{uses_ai}, \code{provides_ai}, \code{processes_personal_data}, \code{is_public_sector} & Trijų būsenų požymiai: tuščia reiškia, kad niekas nepasakė, ir tada atranka palieka aktus, kuriuos požymis lemtų, užuot tylą palaikiusi \enquote{ne}. \code{ai_description} aprašo, kaip naudojamas DI. \\
\code{financial_market} & JSON: finansų rinkos klasifikacija (rinka, paslaugų kategorijos, leidimas). \\
\code{attributes} & JSON: atviras pokalbio metu sužinotų faktų, kuriems dar nėra stulpelio, rinkinys. \\
\end{xltabular}

\paragraph{Ką įmonė turi.} \code{company_credentials} eilutė yra sertifikatas, licencija,
leidimas ar registracija (\code{kind}: \code{certificate}, \code{license}, \code{permit},
\code{registration}) su išdavėju, identifikatoriumi, galiojimu ir apimtimi. Jos
\code{status} yra \code{claimed}, kol žmogus nepatikrina jo išdavėjo registre ir nenustato
\code{verified} (su \code{verified_by}, \code{verified_at}); deklaruotas leidimas laikomas
įrodymu tik su šia išlyga, o \code{evidence_document} gali rodyti į jį įrodantį dokumentą.
\code{company_standards} eilutė yra teisės aktų bazės standartas, kurį įmonė taiko, po vieną
įmonei ir standartui (unikalu), su \code{status} \code{certified}, \code{implementing} arba
\code{planned}, jo taikymo sritimi ir jį sertifikuojančiu leidimu; ataskaita, kuri
neįvardija savų aktų, vertina kiekvieną taikomą standartą, o ataskaita, kuri savo aktus
įvardija, vertina tarp jų esančius standartus, taikomus ar ne (\code{adopted_standards} faile
\file{analysis/pipeline.py}). \code{company_roles} eilutė yra rolės profilis --- pavyzdžiui,
atitikties pareigūno ---, sutrauktas į \code{context}, tekstą, kurį gali perimti pokalbis, o
\code{sources} išvardija dokumentų skirsnius, iš kurių jis sudarytas
(\code{[{documentId, title, section}]}); pavadinimas įmonėje unikalus.

\begin{xltabular}{\textwidth}{@{}L{38mm}Y@{}}
\caption{\code{company_documents}: vienas įkeltas failas}\label{tab:data:cd}\\
\toprule
\hd{Stulpelis} & \hd{Reikšmė} \\
\midrule
\endfirsthead
\toprule
\hd{Stulpelis} & \hd{Reikšmė} \\
\midrule
\endhead
\bottomrule
\endlastfoot
\code{owner}, \code{company} & Jį įkėlęs naudotojas (\code{PROTECT}) ir įmonė, kuriai jis priskirtas (tuščia, kol nepriskirtas). \\
\code{title}, \code{kind} & \code{kind}: \code{policy}, \code{procedure}, \code{contract}, \code{register}, \code{other}. \\
\code{file_path}, \code{mime}, \code{size}, \code{sha256} & Failas įkėlimų aplanke ir jo maiša. \\
\code{text}, \code{char_count} & Ištrauktas tekstas --- tai, ką skaito modelis ir kas skaidoma į fragmentus. \\
\code{extraction_status}, \code{extraction_error} & \code{pending}, \code{done}, \code{empty}, \code{failed}. \\
\code{inventory}, \code{inventory_key} & JSON: įrodymų inventorius ir raktas, kuriam jis apskaičiuotas, \code{<sha256 of the text>:<prompt version id>:s<SCHEMA_VERSION>}; naujas tekstas, instrukcijos versija ar inventoriaus struktūra jį pasendina (\file{companies/inventory.py}, \code{SCHEMA_VERSION = 2}). \\
\code{requirements_key} & Paskutinio jo vidaus reikalavimų perskaitymo \code{<sha256 of the file>:<prompt version id>}; tuščia reiškia, kad niekada neperskaityta. \\
\code{indexed_at} & Kada paskutinį kartą sudaryti jo fragmentai. \\
\code{outline} & JSON: tekste rastos antraštės, \code{[{title, level, offset}]}, kad skaitytojas galėtų peršokti į \enquote{§7.2}. \\
\code{uploaded_at} & Įkėlimo momentas; kartu su \code{owner} --- savininko sąrašo tvarka (indeksuota). \\
\end{xltabular}

Inventorius yra išsaugotas instrukcijos \prompt{report.inventory} atsakymas; jo raktus nustato
šios instrukcijos JSON schema:

\begin{codeblock}
inventory = {documentType: policy | procedure | register | contract | plan
                           | report | profile | other,
             summary, scope, controls: [{name, description, quote}],
             roles, registers, approval, reviewCycle, lawsReferenced,
             citations: [{act, article, paragraph, point, claim, quote}],
             weaknesses}
\end{codeblock}

\code{roles}, \code{registers}, \code{lawsReferenced} ir \code{weaknesses} yra trumpų tekstų
sąrašai; \code{citations} --- nuorodos, kurias tikrina teisės aktų nuorodų patikra
(\cref{ch:company}). \code{company_document_chunks} eilutė yra teksto langas su savo
\code{ordinal}, \code{char_start}, artimiausia virš jo esančia antrašte \code{heading}
(skirsniu, kurį įvardija nuoroda į fragmentą), \code{text} ir \code{embedding}, po jau
aprašytais HNSW ir GIN indeksais bei indeksu pagal dokumentą ir eilės numerį.

\section{Ataskaitos ir veiksmų registras}\label{sec:data:analysis}

Programėlė \code{analysis} laiko ataskaitas ir viską, iš ko ataskaita sudaryta, taip pat
veiksmų registrą, kuris jas pergyvena. Viena ataskaita yra viena \code{reports} eilutė su jos
aktais (\code{report_acts}), išvadomis (\code{report_findings}), veiksmais
(\code{report_actions}) ir nuorodų patikros eilutėmis (\code{report_citations});
\code{company_actions} priklauso įmonei, ne ataskaitai. Kaip ataskaita rengiama, aprašyta
\cref{ch:report}, kaip toliau gyvuoja registras --- \cref{ch:register}.

\begin{xltabular}{\textwidth}{@{}L{38mm}Y@{}}
\caption{\code{reports}: viena įmonės atitikties spragų analizė}\label{tab:data:re}\\
\toprule
\hd{Stulpelis} & \hd{Reikšmė} \\
\midrule
\endfirsthead
\toprule
\hd{Stulpelis} & \hd{Reikšmė} \\
\midrule
\endhead
\bottomrule
\endlastfoot
\code{company}, \code{created_by}, \code{job} & Įmonė, ją pradėjęs naudotojas ir paskutinio jos vykdymo užduotis (\cref{ch:arch}); užduotis pakeičiama, kai į eilę įtraukiama antroji vykdymo dalis. \\
\code{status} & \code{queued}, \code{running}, \code{scope_review} (laukiama, kol žmogus patvirtins aktų sąrašą), \code{done}, \code{failed}, \code{cancelled}. \\
\code{language} & \code{lt} arba \code{en}: kalba, kuria rašoma ataskaita. \\
\code{scope} & JSON: ką vykdymui pavesta atlikti (žr. toliau). \\
\code{profile_snapshot} & JSON: įmonė tokia, kokią ją perskaitė vykdymas (žr. toliau). \\
\code{title}, \code{summary} & Pavadinimas ir santrauka (Markdown), kuriuos parašė santraukos instrukcija. \\
\code{stats} & JSON: kiekiai, kaina ir kilmė (žr. toliau). \\
\code{error} & Kodėl vykdymas nepavyko; \code{AI_BUDGET_EXHAUSTED: …}, kai DI tinklų sąsaja atsisakė dėl savo biudžeto. \\
\code{scope_approved_by}, \code{scope_approved_at} & Kas ir kada patvirtino aktų sąrašą. Vykdymą, pradėtą be peržiūros, pažymi pati apdorojimo grandinė: laikas yra, žmogaus nėra. \\
\code{approved_by}, \code{approved_at}, \code{approval_note} & Patvirtinimas; jį vėl panaikina bet kuri vėlesnė išvados peržiūra (pakeista jos būsena, pastaba ar priežastis). \\
\code{created_at}, \code{started_at}, \code{finished_at} & Trys jos momentai; indeksuota pagal įmonę ir sukūrimo laiką. \\
\end{xltabular}

\begin{codeblock}
scope = {documentIds: [ids of legal_documents named by the user],
         review: true,     only when the scope review was asked for
         internal: true}   only when own requirements are to be assessed

stats = {acts: {applies, uncertain, not_applies,
                assessed, upcoming, unconfirmed, indirect},
         findings: {met, partial, missing, unknown, not_applicable},
         openBySeverity: {critical, high, medium, low},
         upcoming: {findings, openBySeverity},  only if an upcoming act has findings
         internal: {findings, openBySeverity},  only if own requirements have findings
         changes: {since, sinceDate, closed: {count, titles}, new, back},
                                                 only if an earlier report exists
         actions, reviewedObligations, tokens, costUsd,
         citations: {checked, match, mismatch, notFound, notInCorpus, unknown},
         provenance: {prompts: [{slug, version, model}], models,
                      corpus: [{act, versionNo, versionDate}]}}
\end{codeblock}

Apimties \code{documentIds} yra teisės aktų bazės aktai, ne įmonės dokumentai: ten įvardytas
aktas vertinamas net tada, kai jo verdiktas \code{uncertain} arba kai įmonei jis privalomas
tik netiesiogiai (\code{assessment_of} faile \file{analysis/pipeline.py}); aktas, kurį atmeta
katalogo atranka arba kurį modelis pripažįsta
netaikomu, nevertinamas. \code{stats} apskaičiuoja \code{report_stats} faile
\file{analysis/pipeline.py}, kai vykdymas baigiasi, ir dar kartą anksčiau, kai vykdymas
sustoja apimties peržiūrai. Išvados skaičiuojamos pagal \emph{faktinę} būseną --- viršų ima
peržiūrėtojo sprendimas --- įmonei galiojančių aktų atžvilgiu; aktų, vertinamų pasiruošimui,
ir įmonės vidaus reikalavimų išvados skaičiuojamos atskirai, po \code{upcoming} ir
\code{internal}. Kiekiai perskaičiuojami, kai peržiūrėtojas pakeičia išvados būseną
(\code{finding_counts}); visa kita --- ne. \code{tokens} ir \code{costUsd} sumuoja visų
užduoties etapų, priskirtų ataskaitai (\code{subject = report:<id>}), modelio iškvietimus,
todėl įskaičiuojamos abi vykdymo su apimties peržiūra dalys. \code{provenance} įvardija
kiekvieną veikusią instrukcijos versiją ir modelį, o kiekvienam aktui, kurį ataskaita vertino,
užuot atmetusi atrankoje, --- perskaitytos redakcijos numerį ir datą. \code{changes}
(\code{changes} faile \file{analysis/register.py}) kiekvieną pareigą lygina su paskutine
ankstesne baigta ataskaita, kuri ją vertino: spragos \code{closed} (uždarytos), spragos
\code{new} (naujos) ir pareigos, kurios iš atitinkančių vėl tapo spragomis (\code{back}),
kiekviena su kiekiu ir ne daugiau kaip penkiais pavadinimais.

\code{profile_snapshot} yra įmonė tokia, kokią ją su visais priedais pateikia API
(\code{serialize_company(…, with_related=True)} faile \file{companies/serializers.py}):
profilio stulpeliai jų API pavadinimais (\code{institutionType}, \code{usesAi},
\code{financialMarket}, \code{attributes}, \dots) ir sąrašai \code{credentials},
\code{standards}, \code{roles} bei \code{documents} --- pastarajame kiekvieno dokumento
metaduomenys, jo inventoriaus \code{documentType} ir vienos eilutės \code{inventorySummary},
niekada ne jo tekstas. Jį įrašo \code{freeze_profile} faile \file{analysis/pipeline.py}:
pirmosios vykdymo dalies etape \code{profile} ir dar kartą, kai prasideda antroji dalis. Antroji
dalis įmonę perskaito iš naujo, todėl ataskaita, parengta su apimties peržiūra, rodo tą įmonę,
pagal kurią vertintos jos išvados, --- su leidimu, patikrintu ataskaitai laukiant, ar tuo metu
pridėtu dokumentu. Ataskaita, vykdoma be peržiūros, jį įrašo vieną kartą.

\begin{xltabular}{\textwidth}{@{}L{38mm}Y@{}}
\caption{\code{report_acts}: vienas aktas taip, kaip jį traktavo ataskaita}\label{tab:data:ra}\\
\toprule
\hd{Stulpelis} & \hd{Reikšmė} \\
\midrule
\endfirsthead
\toprule
\hd{Stulpelis} & \hd{Reikšmė} \\
\midrule
\endhead
\bottomrule
\endlastfoot
\code{document}, \code{version} & Nuorodos į aktą ir perskaitytą redakciją (laikomos tik naršymui). \\
\code{document_title}, \code{short_title}, \code{jurisdiction}, \code{doc_type}, \code{source_url} & Kopijos iš katalogo ataskaitos vykdymo metu. \code{doc_type} \code{standard} žymi taikomą standartą (taikomas todėl, kad įmonė jį taiko, be modelio iškvietimo); \code{internal} --- eilutę, prie kurios kabinami įmonės vidaus reikalavimai. \\
\code{applicability}, \code{confidence}, \code{reasoning}, \code{conditions_to_verify} & Modelio verdiktas: \code{applies}, \code{not_applies} arba \code{uncertain}; pasitikėjimas, apribotas intervalu 0–1; jo pagrindimas; tikslintini faktai. \\
\code{applies_from} & Data, nuo kurios aktas įmonei privalomas, jei ji vėlesnė už vykdymo dieną. \\
\code{assessment} & Ką ataskaita padarė su aktu: \code{assessed}, \code{upcoming}, \code{unconfirmed}, \code{indirect}, \code{none}. Apskaičiuoja kodas iš verdikto; apimties peržiūra gali jį pakeisti, verdikto --- niekada. \\
\code{citations} & JSON: taikymo srities nuostatos, kuriomis remiasi verdiktas, \code{[{unitId, path, documentTitle}]}. \\
\code{screened_out}, \code{order} & Teigiama reikšmė aktui, kurį katalogo atranka atmetė be modelio iškvietimo; rodymo tvarka. \\
\end{xltabular}

\begin{xltabular}{\textwidth}{@{}L{38mm}Y@{}}
\caption{\code{report_findings}: vienos pareigos įvertinimas}\label{tab:data:rf}\\
\toprule
\hd{Stulpelis} & \hd{Reikšmė} \\
\midrule
\endfirsthead
\toprule
\hd{Stulpelis} & \hd{Reikšmė} \\
\midrule
\endhead
\bottomrule
\endlastfoot
\code{act}, \code{obligation}, \code{source_kind} & Nuorodos į ataskaitos akto eilutę ir į pareigą; \code{legal}, \code{standard} arba \code{internal}. \\
\code{document_title}, \code{short_title}, \code{obligation_title}, \code{obligation_text}, \code{category}, \code{deadline} & Akto pavadinimų ir pareigos kopijos (ataskaitos kalba, jei pareiga turi tekstą anglų kalba). \\
\code{obligation_reviewed} & Ar ataskaitos vykdymo metu žmogus buvo patvirtinęs pareigą. \\
\code{status} & Modelio verdiktas: \code{met}, \code{partial}, \code{missing}, \code{unknown} (neįrodyta), \code{not_applicable}. \\
\code{severity} & Nukopijuotas pareigos reikšmingumas. \\
\code{analysis}, \code{recommendation} & Modelio vertinimas ir rekomendacija. \\
\code{evidence}, \code{citations}, \code{consequence} & JSON (žr. toliau). \\
\code{override_status}, \code{reviewer_note}, \code{reviewed_by}, \code{reviewed_at} & Peržiūrėtojo nustatyta būsena ir pastaba; jos viršų ima kiekviename rodinyje ir eksporte. \\
\code{cause} & Kodėl spraga atsirado, peržiūrėtojo žodžiais; rašo tik žmogus. \\
\code{order} & Pirmiausia neatitinkančios, tada iš dalies, neįrodytos, atitinkančios, netaikomos; kiekvienoje grupėje pirmiausia reikšmingesnės; įmonės vidaus reikalavimai po visų kitų (indeksuota kartu su \code{report}). \\
\end{xltabular}

\begin{codeblock}
evidence    = [{kind: "passage",    id: document id,   title, quote, start}
               {kind: "credential", id: credential id, title, quote}
               {kind: "profile",    id: company id,    title, quote}]
citations   = the obligation's citations, with documentTitle and url added
              (legal, standard); copied as they are (internal)
consequence = [{act, path, where, condition, maxFineEur, maxFinePercent,
                benefitMultiple, persons, measures, measuresFrom: [paths],
                via: {act, path}}]      (via: under a level 2 act only)
\end{codeblock}

Citata sąraše \code{evidence} sutrumpinama iki 600 simbolių, o \code{start} yra simbolis, nuo kurio
ji prasideda dokumento tekste, arba tuščia, kai dokumente ji pažodžiui neužrašyta: tada
skaitytojas mato, ką teigė modelis, ir kad citatos nepavyko rasti tekste. Ataskaita,
parengta tuo metu, kai ištraukos buvo indeksuoti fragmentai, vietoj to išsaugojo fragmento
eilės numerį (\code{chunk}), ir API jį išsprendžia pateikdama ataskaitą
(\file{analysis/serializers.py}). \code{consequence} laiko sankcijų pakopas, kurios apima
nuostatą, tokias, kokias teisės aktų bazė turėjo ataskaitos vykdymo metu
(\code{Consequences.of} faile \file{corpus/sanctions.py}): pirmiausia paties akto, tada vieno
kito akto. Jis tuščias aktui, kuris įmonei dar neprivalomas, standarto ar pačios įmonės
reikalavimui, aktui, kurio sankcijos niekada nebuvo perskaitytos, ir nuostatai, kurios
nepasiekia jokia pakopa. 2~lygio akto atveju įrašai yra pagrindinio akto straipsnio, pagal
kurį tas aktas priimtas, įrašai, kiekvienas su \code{via} --- tuo aktu ir straipsniu
(\cref{ch:obligations}); iki 2026~m. lapkričio 16~d. parengta ataskaita tokiam aktui jų
neturi.

\paragraph{Veiksmai ir nuorodų patikros eilutės.} \code{report_actions} eilutė yra vienas
plano žingsnis (\cref{tab:data:rx}); \code{report_citations} eilutė (\cref{tab:data:rc}) yra
viena įmonės dokumento nuoroda į nuostatą su patikros verdiktu. Nuoroda, kurią daro keli
dokumentai, įvertinama vieną kartą, bet kiekvienas dokumentas vis tiek turi savo eilutę.

\begin{xltabular}{\textwidth}{@{}L{38mm}Y@{}}
\caption{\code{report_actions}: vienas veiksmų plano žingsnis}\label{tab:data:rx}\\
\toprule
\hd{Stulpelis} & \hd{Reikšmė} \\
\midrule
\endfirsthead
\toprule
\hd{Stulpelis} & \hd{Reikšmė} \\
\midrule
\endhead
\bottomrule
\endlastfoot
\code{title}, \code{description}, \code{owner_role}, \code{category} & Ką daryti ir kuri rolė už tai atsakinga. \\
\code{priority} & Jo vieta plane, apskaičiuota kodo --- pagal tai, ką jis uždaro (šiandienos spragos pirmiau už pasiruošimą dar netaikomiems aktams, tada pagal reikšmingumą ir būseną), tada pagal jo laikotarpį ---, niekada neklausiant modelio. \\
\code{effort}, \code{timeframe} & \code{low}, \code{medium}, \code{high}; \code{immediate}, \code{30_days}, \code{90_days}, \code{12_months}. \\
\code{finding_ids} & JSON sąrašas išvadų, kurias jis uždaro, identifikatorių; veiksmas, kuris neuždaro jokios išvados (nuorodų taisymas, faktų patvirtinimas), turi tuščią sąrašą. \\
\code{order} & Rodymo tvarka. \\
\end{xltabular}

\begin{xltabular}{\textwidth}{@{}L{38mm}Y@{}}
\caption{\code{report_citations}: viena dokumento nuoroda ir patikros verdiktas}\label{tab:data:rc}\\
\toprule
\hd{Stulpelis} & \hd{Reikšmė} \\
\midrule
\endfirsthead
\toprule
\hd{Stulpelis} & \hd{Reikšmė} \\
\midrule
\endhead
\bottomrule
\endlastfoot
\code{company_document}, \code{document_title} & Dokumentas, kuriame yra nuoroda (saitas ir kopija). \\
\code{act}, \code{article}, \code{paragraph}, \code{point}, \code{claim}, \code{quote} & Nuoroda tiksliai taip, kaip ją užrašė dokumentas, ir ką, anot jo, nuostata apima. \\
\code{status} & \code{found}, \code{not_found}, \code{act_not_in_corpus}. \\
\code{heading_match} & Rastai nuostatai: \code{match}, \code{mismatch}, \code{unknown}; kitais atvejais tuščia. \\
\code{judged_by} & \code{model} arba \code{overlap}, kai vertintojas buvo nepasiekiamas ir sprendė žodžių sutapimas; tuščia, kai nieko nelyginta. \\
\code{legal_document}, \code{unit} & Nuorodos į aktą ir rastą nuostatą. \\
\code{short_title}, \code{unit_path}, \code{heading}, \code{note}, \code{excerpt}, \code{citation_url} & Rastos nuostatos kopijos, vertintojo pastaba, nuostatos ištrauka ir oficiali jos nuoroda. \\
\code{order} & Pirmiausia problemos (neatitinka, nerasta, neaišku, akto nėra bazėje, atitinka), tada dokumentų tvarka. \\
\end{xltabular}

\begin{xltabular}{\textwidth}{@{}L{38mm}Y@{}}
\caption{\code{company_actions}: vienas įmonės registro veiksmas}\label{tab:data:ca}\\
\toprule
\hd{Stulpelis} & \hd{Reikšmė} \\
\midrule
\endfirsthead
\toprule
\hd{Stulpelis} & \hd{Reikšmė} \\
\midrule
\endhead
\bottomrule
\endlastfoot
\code{company}, \code{report}, \code{plan_action_id} & Įmonė; patvirtinta ataskaita, kurios planas jį įtraukė; to plano veiksmo identifikatorius. Unikalu pagal įmonę ir \code{plan_action_id}, todėl dukart įtraukus planą nieko nepridedama. \\
\code{title}, \code{description}, \code{category}, \code{priority}, \code{effort} & Nukopijuota iš plano veiksmo. \\
\code{owner}, \code{due_date} & Kas jį vykdo (iš pradžių --- plane nurodyta atsakinga rolė) ir iki kada (patvirtinimo data plius 14, 30, 90 arba 365 dienos keturiems laikotarpiams); žmogus gali pakeisti abu. \\
\code{status} & \code{open}, \code{in_progress}, \code{done} (atsakingo asmens žodis), \code{closed}. Būsena \enquote{Vėluojantis} nesaugoma: tai atviras arba vykdomas veiksmas, kurio terminas praėjo. \\
\code{items} & JSON: atviros išvados, kurias jis turi uždaryti (žr. toliau). \\
\code{note}, \code{closed_at}, \code{closed_by}, \code{closed_from}, \code{confirmed_report} & Kaip jis uždarytas: žmogaus, nurodžiusio priežastį \code{note}, arba vėlesnės patvirtintos ataskaitos (\code{confirmed_report}); \code{closed_from} --- būsena, iš kurios uždarytas, kad atšaukus patvirtinimą jį būtų galima grąžinti. \\
\code{done_at} & Kada atsakingas asmuo paskutinį kartą pažymėjo jį atliktu: tik po to parengta ataskaita gali jį grąžinti vykdyti. \\
\end{xltabular}

\begin{codeblock}
items = [{obligationId, findingId, title, act, path, sourceKind, status,
          confirmed: null | {reportId, at}}]
\end{codeblock}

\section{Programos dokumentai}\label{sec:data:programme}

Programėlė \code{programme} turi vieną lentelę. \code{work_products} eilutė yra viena įmonės
programos dokumento versija --- atitikties politikos, valdymo struktūros su atitikties procesu
ir RACI arba rizikos registro (\cref{ch:programme}). Jos \code{kind} yra \code{policy},
\code{governance} arba \code{risks}, be apribojimo; \code{version} didėja kiekvienai įmonei ir
rūšiai (kartu unikalūs); \code{language} yra \code{lt} arba \code{en}; \code{status} yra
\code{queued}, \code{running}, \code{draft} (sugeneruotas ar redaguojamas, laukia žmogaus),
\code{approved} (užfiksuotas: pakeitimas yra nauja versija), \code{failed} arba
\code{cancelled}. Be \code{job}, \code{title} ir \code{error}, ji turi keturis JSON stulpelius
ir tris žymas, kas ką padarė: \code{created_by}, \code{edited_by} su \code{edited_at} ir
\code{approved_by} su \code{approved_at} bei \code{approval_note}.

\begin{codeblock}
content (policy)     = {title,
                        parts: [{key, heading, text, decisions, sources,
                                 legalBasis}],
                        hierarchy: [{documentId, document, level, topic, note}],
                        missing: [{level, title, purpose, supports, priority,
                                   legalBasis, partlyCoveredBy}]}
content (governance) = {title, positions, gaps, structure, process, raci}
content (risks)      = {title, criteria: {likelihood, impact, calibration},
                        appetite, scenarios}
checks = [{rule, ok, blocking, detail}]
basis  = {reportId, products: {<kind>: {id, version, status}},
          documents: [{id, title, sha256}]}
stats  = {tokens, costUsd, provenance: {prompts: [{slug, version, model}],
                                        models}}
\end{codeblock}

\code{content} yra dokumentas tokios išsaugomos struktūros, kokią jam suteikia jo rūšies
\code{validate} (\file{programme/policy.py}, \file{governance.py}, \file{risks.py});
\code{checks} --- rūšies taisyklės taip, kaip turinys jas paskutinį kartą atitiko ar pažeidė,
o pažeista \code{blocking} taisyklė neleidžia dokumento patvirtinti; \code{basis} įvardija, ką
vykdymas skaitė, --- ataskaitą, ankstesnius programos dokumentus, įmonės dokumentus su jų maišomis ---, kad
skaitytojas matytų, kuo dokumentas remiasi.

\section{Lentelės, kuriomis sistema veikia}\label{sec:data:ops}

Likusios lentelės, pavaizduotos \cref{fig:data:ops}, yra paskyros ir sesijos, instrukcijos,
eilė, modelio iškvietimų įrašai, pokalbiai ir veiksmų žurnalas.


\paragraph{Naudotojai ir sesijos.} \code{users} eilutė yra paskyra, ir kiekviena kita
lentelė, kuri įvardija asmenį, rodo į ją.

\begin{xltabular}{\textwidth}{@{}L{38mm}Y@{}}
\caption{\code{users} ir \code{sessions}}\label{tab:data:us}\\
\toprule
\hd{Stulpelis} & \hd{Reikšmė} \\
\midrule
\endfirsthead
\toprule
\hd{Stulpelis} & \hd{Reikšmė} \\
\midrule
\endhead
\bottomrule
\endlastfoot
\multicolumn{2}{@{}l}{\textit{users}} \\
\code{email}, \code{display_name} & El. pašto adresas unikalus, ir unikalus nepaisant raidžių dydžio (funkcinis indeksas pagal jo mažąsias raides). \\
\code{password_hash} & Slaptažodžio maiša; niekada ne pats slaptažodis. \\
\code{role} & \code{admin} arba \code{user}. \\
\code{active} & Išjungiklis, kurį iš naujo tikrina kiekviena autentifikuota užklausa; paprastas paskyros trynimas ją išjungia, ir tik administratoriaus galutinis ištrynimas ją ištrina su jos duomenimis (\cref{tab:data:ondelete}). \\
\code{daily_budget_usd} & Kiek paskyra per UTC parą gali išleisti modelio iškvietimams: tuščia reiškia diegimo numatytąją reikšmę (\code{LLM_DAILY_BUDGET_USD}), 0 --- neribota (\cref{ch:models}). \\
\code{created_at}, \code{updated_at}, \code{last_login_at} & \\
\midrule
\multicolumn{2}{@{}l}{\textit{sessions}} \\
\code{user}, \code{token_hash} & Paskyra ir kliento turimo prieigos rakto SHA-256 (unikalus); pats raktas niekada nesaugomas. \\
\code{expires_at}, \code{created_at}, \code{user_agent}, \code{ip} & Galiojimo pabaiga ir iš kur atėjo sesija; indeksuota pagal naudotoją ir galiojimo pabaigą. Prisijungiant išlaikoma tik \code{SESSIONS_PER_USER = 10} naujausių paskyros sesijų (\file{users/auth.py}). \\
\end{xltabular}

\paragraph{Instrukcijos.} Nurodymai, kuriuos gauna modelis, saugomi kaip instrukcijos su
nekeičiamomis versijomis (\cref{tab:data:pr}); redaguojant instrukciją pridedama versija, o
aktyvinant versiją perkeliamas \code{active_version}. Kaip pristatomos numatytosios
instrukcijos tampa versijomis, aprašyta \cref{ch:prompts}. 2026 m. lapkričio 13 d. lentelėse
buvo 23 instrukcijos ir 85 versijos, ir kiekviena aktyvi versija buvo pristatyta numatytoji instrukcija;
nauja numatytoji pokalbio instrukcijos versija buvo įkelta 2026~m. lapkričio 16~d. ir dar kartą
lapkričio 17~d., naujos numatytosios sankcijų instrukcijos (dukart) ir inventoriaus
instrukcijos versijos --- lapkričio 19~d., o dar viena sankcijų instrukcijos versija ---
lapkričio 20~d., todėl jose yra 91 versija, ir kiekviena aktyvi
versija tebėra pristatyta
numatytoji instrukcija.

\begin{xltabular}{\textwidth}{@{}L{38mm}Y@{}}
\caption{\code{prompts} ir \code{prompt_versions}}\label{tab:data:pr}\\
\toprule
\hd{Stulpelis} & \hd{Reikšmė} \\
\midrule
\endfirsthead
\toprule
\hd{Stulpelis} & \hd{Reikšmė} \\
\midrule
\endhead
\bottomrule
\endlastfoot
\multicolumn{2}{@{}l}{\textit{prompts}} \\
\code{slug}, \code{name}, \code{description} & Unikalus pavadinimas, kuriuo ją kviečia kodas (\prompt{report.gap_analysis}). \\
\code{kind} & \code{chat}, \code{pipeline} arba \code{corpus}. \\
\code{variables}, \code{output_schema} & JSON: kintamieji, kuriuos ji gali naudoti, ir JSON schema, kurią turi atitikti jos atsakymas. \\
\code{active_version} & Naudojama versija. \\
\midrule
\multicolumn{2}{@{}l}{\textit{prompt\_versions}} \\
\code{prompt}, \code{version_no} & Kartu unikalūs. \\
\code{text}, \code{model}, \code{temperature}, \code{max_tokens} & Tekstas ir modelio parametrai; \code{model} yra \code{provider:model} formos eilutė. \\
\code{source}, \code{default_hash} & \code{default} --- tekstui, pristatytam faile \file{prompts/defaults/*.yaml} (maiša įvardija tikslų pristatytą tekstą); \code{user} --- redaktoriuje parašytam tekstui. \\
\code{note}, \code{created_by}, \code{created_at} & \\
\end{xltabular}

\paragraph{Užduotys ir etapai.} Eilę sudaro dvi lentelės, o \cref{ch:arch} aprašo, kaip jos
veikia kaip viena: užduoties paėmimas, užrakinimas, gyvybės signalai, pakartojimai ir
atšaukimas. \Cref{tab:data:jobs} išvardija jų stulpelius; tarp jų esantys JSON stulpeliai yra
tokios struktūros:

\begin{codeblock}
jobs.progress   = {step, percent, message}
jobs.result     = what the handler returned; report.analyze:
                  {reportId, findings, actions} or {reportId, waiting: "scope"};
                  programme.generate: {productId, checks: how many broken}
job_steps.input_ref, job_steps.output
                = what the step was given and what it counted; the
                  gap_analysis output, for instance: {met, partial, missing,
                  unknown, not_applicable, documents: "all" | "nearest",
                  quotesPlaced, quotesUnplaced}
job_steps.usage = {calls, inputTokens, outputTokens, costUsd,
                   models: [{model, calls, inputTokens, outputTokens, costUsd}]}
\end{codeblock}

\begin{xltabular}{\textwidth}{@{}L{38mm}Y@{}}
\caption{\code{jobs} ir \code{job_steps}}\label{tab:data:jobs}\\
\toprule
\hd{Stulpelis} & \hd{Reikšmė} \\
\midrule
\endfirsthead
\toprule
\hd{Stulpelis} & \hd{Reikšmė} \\
\midrule
\endhead
\bottomrule
\endlastfoot
\multicolumn{2}{@{}l}{\textit{jobs}} \\
\code{kind}, \code{payload}, \code{subject} & Apdorojimo funkcija, jos JSON įvestis ir tai, ką užduotis liečia (\code{report:<id>}; indeksuota). \\
\code{status} & \code{queued}, \code{running}, \code{done}, \code{failed}, \code{cancelled}. \\
\code{priority}, \code{run_after} & Paėmimo tvarka ir anksčiausias momentas; paėmimą aptarnauja vienas indeksas (būsena, \code{run_after}, prioritetas, sukūrimas). \\
\code{attempts}, \code{max_attempts}, \code{cancel_requested}, \code{locked_by}, \code{heartbeat_at} & Pakartojimai, atšaukimas ir ją laikančio užduočių vykdytojo užraktas. \\
\code{progress}, \code{result}, \code{error} & JSON eiga ir rezultatas; klaidos tekstas. \\
\code{created_by}, \code{created_at}, \code{started_at}, \code{finished_at} & Kas ją sukėlė (tuščia planuoklio užduotims) ir jos momentai. \\
\midrule
\multicolumn{2}{@{}l}{\textit{job\_steps}} \\
\code{job}, \code{ordinal}, \code{name} & Etapo vieta užduotyje (indeksuota): ataskaitai \code{profile}, \code{documents}, \code{citations}, \code{applicability}, … \\
\code{status} & \code{running}, \code{done}, \code{failed}, \code{skipped}. \\
\code{prompt_version_id}, \code{model} & Jo iškvietimų instrukcijos versija ir modelis, tekstu: etapo įrašas pergyvena versijos eilutę. \\
\code{input_ref}, \code{output}, \code{usage} & JSON, kaip nurodyta aukščiau; \code{usage} etapui pasibaigus nuskaitomas iš jo \code{llm_calls} eilučių. \\
\code{duration_ms}, \code{error}, \code{started_at}, \code{finished_at} & \\
\end{xltabular}

\paragraph{Modelio iškvietimų įrašai.} \code{llm_calls} eilutė yra vienas modelio
iškvietimas, ir iš šios lentelės skaičiuojamas sąnaudų puslapis, dienos biudžetas ir
kiekvienas šios knygos kainos skaičius (\cref{ch:models,ch:control}).

\begin{xltabular}{\textwidth}{@{}L{38mm}Y@{}}
\caption{\code{llm_calls}: vienas modelio iškvietimas}\label{tab:data:lc}\\
\toprule
\hd{Stulpelis} & \hd{Reikšmė} \\
\midrule
\endfirsthead
\toprule
\hd{Stulpelis} & \hd{Reikšmė} \\
\midrule
\endhead
\bottomrule
\endlastfoot
\code{kind}, \code{purpose} & \code{chat}, \code{job_step}, \code{embed}, \code{prompt_test}; trumpa laisva žyma (ne daugiau kaip 80 simbolių). \\
\code{user} & Apmokestinama paskyra; tuščia darbui, kuris niekam nepriskiriamas, pavyzdžiui, teisės aktų bazės pareigų išskyrimui ataskaitos viduje. \\
\code{thread_id}, \code{job_step_id}, \code{prompt_version_id} & Kam skirtas iškvietimas, tekstiniais identifikatoriais be išorinio rakto; \code{job_step_id} indeksuotas, nes per jį nuskaitomos etapo sąnaudos ir ataskaitos kaina. \\
\code{model} & Modelis taip, kaip jį vadina DI tinklų sąsaja. \\
\code{input_tokens}, \code{cached_input_tokens}, \code{cache_write_tokens}, \code{output_tokens}, \code{reasoning_tokens} & Žetonų skaičiai; iš įvesties --- iš tiekėjo instrukcijų podėlio nuskaityti ir į jį įrašyti, už kuriuos mokama kitomis kainomis. \\
\code{cost_usd} & Kaina JAV doleriais, šeši skaitmenys po kablelio, apskaičiuojama pagal DI tinklų sąsajos kainoraštį įrašant eilutę. \\
\code{duration_ms}, \code{tool_calls}, \code{outcome} & Trukmė, JSON sąrašas įrankių, kuriuos iškvietė pokalbio ėjimas, ir rezultatas (numatytasis \code{ok}; laisvas tekstas iki 40 simbolių). \\
\code{transcript} & JSON pokalbio ėjimui: \code{{user, assistant}}, klausimo ir atsakymo tekstas --- įrašas, ką asistentas kam pasakė. Ištrinamas, kai paskyra galutinai ištrinama (\file{users/purge.py}); pati eilutė lieka, atsieta. \\
\code{created_at} & Indeksuota atskirai ir kartu su \code{user}. \\
\end{xltabular}

Šalia stovi dvi mažos tos pačios programėlės lentelės. \code{model_prices}, kurios raktas yra
modelio identifikatorius, laiko paskutinę žinomą kiekvieno modelio kainą JAV doleriais už
žetoną --- įvesties, išvesties ir, kai DI tinklų sąsaja jas nurodo, iš podėlio nuskaitytos ir
į jį įrašytos įvesties ---, atnaujinamą iš sąsajos kainoraščio ir naudojamą, kai kainoraštis
nepasiekiamas. \code{mcp_servers} laiko išorinį MCP serverį, kurio įrankius gali kviesti
pokalbių asistentas: unikalų \code{name} (jo įrankių priešdėlį), \code{url},
\code{transport} (\code{http} arba \code{sse}), \code{headers} su jo prisijungimo duomenimis,
\code{allowed_tools} (tuščia reiškia visus įrankius), \code{enabled} ir paskutinį sėkmingą jo
įrankių sąrašą (\code{tools}, \code{[{name, description}]}).

\paragraph{Pokalbiai.} Pokalbis yra vieno naudotojo gija su jos žinutėmis
(\cref{tab:data:ct}). Konsolė prideda naudotojo žinutes, o asistento paslauga --- savo
atsakymus, abi per \file{chat/store.py}, kuris atmeta žinutę, kurios turinyje nėra lygiai
vienos rolės, ir išsaugotą žinutę, kuri pakeistų savo rolę (\cref{ch:chat}).

\begin{xltabular}{\textwidth}{@{}L{38mm}Y@{}}
\caption{\code{chat_threads} ir \code{chat_messages}}\label{tab:data:ct}\\
\toprule
\hd{Stulpelis} & \hd{Reikšmė} \\
\midrule
\endfirsthead
\toprule
\hd{Stulpelis} & \hd{Reikšmė} \\
\midrule
\endhead
\bottomrule
\endlastfoot
\multicolumn{2}{@{}l}{\textit{chat\_threads}} \\
\code{user} & Vienintelis naudotojas, kuriam priklauso pokalbis. \\
\code{company}, \code{role} & Pasirinktinai --- įmonė, apie kurią jis vyksta, ir rolė (\code{company_roles}), kurios kontekstą jis perima. \\
\code{title} & Pavadinimą suteikia pirmoji žinutė, modelis arba naudotojas. \\
\code{model} & Naudotojo pasirinktas DI tinklų sąsajos modelis (\code{gateway:azure/…}); tuščia reiškia pačios pokalbio instrukcijos modelį. \\
\code{message_count}, \code{archived}, \code{created_at}, \code{updated_at} & Indeksuota pagal naudotoją, archyvavimą ir paskutinį atnaujinimą. \\
\midrule
\multicolumn{2}{@{}l}{\textit{chat\_messages}} \\
\code{id} & Identifikatorius, kurį žinutei suteikė pokalbio klientas (žr. \cref{sec:data:conventions}). \\
\code{thread}, \code{parent_id} & Pokalbis ir žinutė, po kurios ji eina, tekstu. \\
\code{format} & Formato pavadinimas, kurį nurodo rašantysis; konsolė ir asistento paslauga rašo \code{ai-sdk/v6}. \\
\code{role} & Vienintelė žinutės rolė, išsaugant nuskaitoma iš jos turinio. \\
\code{content} & JSON: žinutė tokia, kokią ją laiko pokalbio klientas, ne daugiau kaip \code{MAX_CONTENT_CHARS = 400000} simbolių. \\
\code{created_at} & Išlaikoma, kai žinutė siunčiama dar kartą, kad istorija niekada nepersirikiuotų (indeksuota kartu su pokalbiu). \\
\end{xltabular}

\paragraph{Veiksmų žurnalas ir užklausų ribotuvas.} \code{audit_log} eilutė įrašo vieną
pakeitimą, kurį administratoriui gali prireikti pamatyti, --- prisijungimą, naudotojo
pakeitimą, instrukcijos aktyvinimą, teisės aktų bazės pakeitimą, ataskaitos patvirtinimą
(\cref{tab:data:au}). Ji rašoma tik per \code{audit()} faile \file{audit/log.py}.
\code{rate_limit_windows} eilutė yra vienas fiksuotas ribotuvo langas: jo \code{key}, jame
gautų užklausų skaičius \code{count} ir \code{reset_at}; jį bendrai naudoja visi API
žiniatinklio procesai (\file{common/ratelimit.py}, \cref{ch:security}).

\begin{xltabular}{\textwidth}{@{}L{38mm}Y@{}}
\caption{\code{audit_log}: vienas užfiksuotas veiksmas}\label{tab:data:au}\\
\toprule
\hd{Stulpelis} & \hd{Reikšmė} \\
\midrule
\endfirsthead
\toprule
\hd{Stulpelis} & \hd{Reikšmė} \\
\midrule
\endhead
\bottomrule
\endlastfoot
\code{user} & Kas tai padarė; tuščia veiksmui be prisijungusio naudotojo. \\
\code{action} & Pavadinimas su taškais: \code{reports.approve}, \code{reports.approve_scope}, \code{corpus.delete}, \dots \\
\code{target_type}, \code{target_id} & Su kuo tai padaryta, tekstu (indeksuota kartu). \\
\code{details} & JSON: kas pasikeitė, taip, kaip tai užfiksuoja veiksmas, --- apimties patvirtinimo atveju kiekvienas aktas, kurio traktavimą žmogus pakeitė. \\
\code{ip}, \code{created_at} & Iš kur ir kada (indeksuota, pirmiausia naujausi). \\
\end{xltabular}

\section{Ką ataskaita nusikopijuoja ir ko ne}\label{sec:data:snapshots}

Atitikties ataskaita yra įrodymas, kas rasta tam tikrą dieną, ir vėliau ji turi skaitytis
taip pat, kai teisės aktų bazėje atsiranda nauja akto redakcija, pareiga performuluojama,
dokumentas pakeičiamas ar pasikeičia įmonės profilis. Modelių failai šią taisyklę įtvirtina:
ataskaita \enquote{išsaugo kiekvieno rodomo pavadinimo ir teksto kopiją (aktų pavadinimų,
pareigų tekstų, nuorodų kelių, nuostatų antraščių), todėl ji skaitosi taip pat ir pasikeitus
teisės aktų bazei, pareigoms, dokumentams ar profiliui; išoriniai raktai laikomi tik
naršymui} (\file{analysis/models.py}).

\begin{codebox}{Kas nukopijuojama į ataskaitą jos vykdymo metu}
\begin{itemize}
  \item Įmonė: \code{reports.profile_snapshot}, visas profilis su leidimais, taikomais
    standartais, rolėmis ir dokumentų sąrašu su jų inventorių santraukomis. Ataskaita, sustojusi
    apimties peržiūrai, jį nukopijuoja dar kartą, kai prasideda jos antroji dalis.
  \item Kiekvienas aktas: jo pavadinimas, trumpasis pavadinimas, jurisdikcija, tipas ir
    nuoroda (\code{report_acts}) bei perskaityta redakcija (saitas, o jos numeris ir data ---
    \code{stats.provenance.corpus}).
  \item Kiekviena įvertinta pareiga: jos pavadinimas ir tekstas ataskaitos kalba, kategorija,
    reikšmingumas ir terminas, ar žmogus ją buvo patvirtinęs, ir nuostatos, kurias ji cituoja,
    su jų keliais ir oficialiomis nuorodomis (\code{report_findings}).
  \item Kiek gali kainuoti pažeidimas: sankcijų pakopos, kurios tą dieną apėmė nuostatą
    (\code{report_findings.consequence}).
  \item Įrodymai: kiekviena modelio pateikta citata su vieta, kurioje ji tą dieną buvo savo
    dokumente (\code{start}).
  \item Teisės aktų nuorodų patikra: kiekviena nuoroda taip, kaip užrašyta, rasta nuostata,
    jos antraštė, ištrauka ir nuoroda (\code{report_citations}).
  \item Kuri instrukcijos versija ir modelis atliko kiekvieną etapą ir kiek kainavo vykdymas
    (\code{stats.provenance}, \code{tokens}, \code{costUsd}).
\end{itemize}
\end{codebox}

Tai, kas nenukopijuojama, taip pat apgalvota. Visas akto ir įmonės dokumento tekstas lieka
savo lentelėje: ataskaita laiko citatas, ištraukas ir poslinkius į juos, ne pačius tekstus.
Citatos \code{start} yra poslinkis į dokumento tekstą tokį, koks jis buvo; ištrynus dokumentą,
nuoroda išvaloma, o citata lieka. Instrukcijos įvardijamos, bet nekopijuojamos: kilmė nurodo
trumpąjį pavadinimą, versijos numerį ir modelį, o pati versija išlieka, nes instrukcijų
versijos nekeičiamos ir jokio maršruto, kuris versiją ištrintų, nėra. Užduoties etapai, kuriuose saugomas kiekvieno etapo įrašas, neišvalomi, kol ataskaita
rodo į užduotį (\cref{ch:arch}). Du dalykai pateikiamoje ataskaitoje sąmoningai
apskaičiuojami iš naujo: \code{outdatedActs} --- aktai, kurie gavo naujesnę redakciją po to,
kai ataskaita juos perskaitė (kad skaitytojas žinotų, jog ją reikia vykdyti iš naujo), ir
kiekvienos išvados \code{history} --- kaip jos pareiga atrodė ankstesnėse įmonės ataskaitose
(\file{analysis/serializers.py}).

Veiksmų registras kopijuoja taip pat: \code{company_actions.items} elementas laiko išvados
pavadinimą, aktą ir nuostatos kelią tokius, kokie jie buvo patvirtinimo metu, o pareigą
įvardija identifikatoriumi. Niekada nepatvirtinta pareiga pakeičiama, kai jos aktas
perskaitomas iš naujo, todėl jos identifikatorius gali išnykti; tada API pateikia elementą be
jo pareigos, kaip tokį, kurio jokia vėlesnė ataskaita negali patvirtinti (\code{standing}
faile \file{analysis/register.py}).

\section{Kas rašo kuriuos stulpelius}\label{sec:data:writers}

Lentelėse sumaišytos trys turinio rūšys, o riba tarp jų yra ta pati, kurią ši knyga brėžia
visur: ką parašė modelis, ką apskaičiavo kodas ir ką nusprendė žmogus. Kodas apdorojimo
grandinės instrukcijos atsakymą išsaugo tik tada, kai jis patikrintas pagal griežtąją
instrukcijos JSON schemą, jo reikšmės patikrintos pagal leistinas, o kiekviena modelio
pateikta nuoroda (nuostata \code{P3}, pareiga \code{O7}, dokumentas \code{D4}) susieta su tikra
eilute; nuoroda, kuri su niekuo nesusiejama, atmetama.

\begin{aibox}{Stulpeliai, kuriuose saugoma tai, ką atsakė modelis}
\code{legal_document_versions.scope} ir \code{.sanctions}; išskirtos \code{obligations}
eilutės tekstiniai laukai (pavadinimas, tekstas, adresatas, sąlygos, terminas, įrodymai,
kategorija, reikšmingumas, \code{applies_to}) ir jos \code{verification};
\code{company_documents.inventory}; ataskaitoje --- \code{report_acts.applicability},
\code{confidence}, \code{reasoning}, \code{conditions_to_verify} ir \code{applies_from},
\code{report_findings.status}, \code{analysis}, \code{recommendation} ir \code{evidence}
citatos, \code{report_citations.heading_match} ir \code{note}, kai \code{judged_by} yra
\code{model}, \code{report_actions} tekstas ir tai, kurie parengti veiksmų projektai
sujungiami į vieną (\prompt{report.plan_merge}), bei \code{reports.title} ir \code{summary};
\code{work_products.content}, kol jo neredagavo žmogus; \code{chat_threads.title}, kai
pokalbiui pavadinimą suteikė modelis, ir asistento žinutės.
\end{aibox}

\begin{codebox}{Stulpeliai, kuriuos apskaičiuoja kodas}
\code{report_acts.assessment} (iš verdikto, akto tipo ir ataskaitos apimties);
\code{report_findings.consequence}, \code{order} ir kiekvienos citatos \code{start};
\code{report_actions.priority} ir \code{finding_ids} (sujungtas veiksmas uždaro tai, ką
uždarė jo projektai, o išvada, kuriai planavimo iškvietimas nepaskyrė veiksmo, prijungiama
prie artimiausio, \code{_nearest_action} faile \file{analysis/pipeline.py});
\code{reports.stats}; \code{legal_document_versions.changes}; \code{work_products.checks};
\code{company_actions} elementai, terminai ir jų patvirtinimas vėlesne ataskaita; kiekvienas
kiekis, raktas, maiša ir kaina.
\end{codebox}

\begin{humanbox}{Stulpeliai, kuriuose užfiksuotas žmogaus sprendimas}
\code{obligations.status} (patvirtinimas ar atmetimas), \code{reviewed_by} ir \code{notes};
\code{company_credentials.status} reikšmė \code{verified}; apimties peržiūra
(\code{reports.scope_approved_by} ir jos pakeistų aktų \code{assessment}); išvados
\code{override_status}, \code{reviewer_note} ir \code{cause}; patvirtinimas
(\code{reports.approved_by}, \code{approval_note}); registro veiksmo atsakingas asmuo,
terminas, būsena ir uždarymo pastaba; programos dokumento redagavimas ir patvirtinimas
(\code{edited_by}, \code{approved_by}); šaltinio \code{user} instrukcijos versija. Nė vieno
iš jų nenustato jokia apdorojimo grandinė. Tris iš jų galima nustatyti ir per pokalbį, kai modelis
naudotojo vardu ir su naudotojo teisėmis iškviečia įrankį (\cref{ch:chat}): ataskaitos
apimties patvirtinimas (\code{approveReportScope}, naudotojui priėmus sąrašą ir jo kainą),
išvados būsenos pakeitimas (\code{overrideFinding}, su naudotojo nurodyta priežastimi) ir
leidimo pažymėjimas patikrintu (\code{updateCredential} arba \code{addCredential} su ta
būsena, naudotojui pasakius, kad leidimas buvo patikrintas). Tada eilutėje įvardijamas naudotojas, lyg būtų naudotasi konsole. Joks įrankis
nerašo išvados \code{cause}.
\end{humanbox}

\section{Kaip atsirado schema}\label{sec:data:history}

Lenteles kuria ir keičia Django migracijos, po vieną aplanką kiekvienai programėlei. Jų
pavadinimai yra trumpa modelio istorija (\cref{tab:data:migrations}). 2026 m. lapkričio 13 d.
veikianti duomenų bazė buvo pasiekusi paskutinę čia išvardytą migraciją programėlėse
\code{analysis}, \code{assistant}, \code{corpus}, \code{companies}, \code{programme},
\code{prompts}, \code{users}, \code{chat}, \code{audit}, \code{jobs} ir \code{ops}.

\begin{xltabular}{\textwidth}{@{}L{20mm}Y@{}}
\caption{Kiekvienos programėlės migracijos pagal failų pavadinimus}\label{tab:data:migrations}\\
\toprule
\hd{Programėlė} & \hd{Migracijos po \code{0001_initial}} \\
\midrule
\endfirsthead
\toprule
\hd{Programėlė} & \hd{Migracijos po \code{0001_initial}} \\
\midrule
\endhead
\bottomrule
\endlastfoot
\code{corpus} & 0002 \code{obligations_extracted_at} · 0003 pakeistos šaltinių rūšys · 0004 \code{obligation_embedding_en} · 0005 \code{standards} · 0006 redakcijos \code{scope} · 0007 pareigos \code{verification} · 0008 redakcijos \code{sanctions} · 0009 redakcijos \code{changes} \\
\code{companies} & 0002 dokumento \code{indexed_at}, \code{inventory} ir kt. · 0003 \code{companyrole}, leidimo įrodymų dokumentas ir kt. · 0004 \code{companystandard} · 0005 dokumento \code{requirements_key} \\
\code{analysis} & 0002 \code{initial} · 0003 \code{finding_obligation_reviewed} · 0004 ataskaitos patvirtinimo pastaba ir laikas · 0005 \code{recount_effective_findings} · 0006 \code{reportcitation} · 0007 akto \code{doc_type} · 0008 akto \code{assessment} · 0009 ataskaitos \code{scope_review} · 0010 išvados \code{consequence}, \code{cause} · 0011 \code{companyaction} \\
\code{assistant} & 0002 iškvietimo kaina ir paskirtis · 0003 pakeistos iškvietimų rūšys · 0004 \code{mcpserver} · 0005 \code{modelprice} · 0006 iškvietimų indeksas pagal užduoties etapą · 0007 \code{prompt_cache_accounting} \\
\code{users} & 0002 naudotojo žetonų biudžetas · 0003 žetonų biudžetą pakeitė \code{daily_budget_usd} \\
\code{chat} & 0002 \code{initial} · 0003 pokalbio \code{role} · 0004 pokalbio \code{model} \\
\code{prompts} & 0002 versijos \code{default_hash}, \code{source} \\
\code{programme}, \code{jobs}, \code{audit}, \code{ops} & nėra \\
\end{xltabular}
